\section*{الفصل \ref{ch:g10:lewis-and-shape} --- بنى لويس وأشكال الجزيئات}

\begin{solution}{exo:g10:lewis-and-shape:1}
زوج من \omterm{def:g9:inside-the-atom:nucleus}{الإلكترونات} تتشارك فيه ذرتان. وزوج من \omterm{def:g10:electron-shells:valence}{إلكترونات التكافؤ} يخص \omterm{def:g7:atoms-and-molecules:atom}{ذرة} واحدة
فقط، غير مشترك.
\end{solution}

\begin{solution}{exo:g10:lewis-and-shape:2}
H: 1 (ينقصه \omterm{def:g9:inside-the-atom:nucleus}{إلكترون} واحد لثنائيته). C: 4، وN: 3، وO: 2 (ينقصها 4 و3 و2
\omterm{def:g9:inside-the-atom:nucleus}{إلكترونات} لثمانيتها).
\end{solution}

\begin{solution}{exo:g10:lewis-and-shape:3}
\chemfig{H-\charge{0=\:,90=\:,270=\:}{Cl}}: للكلور الزوج المشترك وثلاثة أزواج
حرة، أي 8 \omterm{def:g9:inside-the-atom:nucleus}{إلكترونات}.
\end{solution}

\begin{solution}{exo:g10:lewis-and-shape:4}
$5 + 3 \times 1 = 8$ \omterm{def:g9:inside-the-atom:nucleus}{إلكترونات}، أي 4 أزواج (ثلاث روابط، \omterm{def:g10:lewis-and-shape:lone-pair}{وزوج حر} واحد).
\end{solution}

\begin{solution}{exo:g10:lewis-and-shape:5}
\ce{CH4}: رباعي الأوجه. \ce{NH3}: هرمي. \ce{H2O}: منحنٍ. \ce{CO2}:
خطي.
\end{solution}

\begin{solution}{exo:g10:lewis-and-shape:6}
\chemfig{H-\charge{90=\:,270=\:}{S}-H}: كما في الماء، أربع مجموعات حول الكبريت
منها زوجان حران: منحنٍ.
\end{solution}

\begin{solution}{exo:g10:lewis-and-shape:7}
$2 \times 5 = 10$ \omterm{def:g9:inside-the-atom:nucleus}{إلكترونات}، أي 5 أزواج. رابطة أحادية و4 أزواج حرة تُبقي لكل
نيتروجين 6 \omterm{def:g9:inside-the-atom:nucleus}{إلكترونات}؛ وتحويل زوجين حرين إلى زوجين مشتركين يعطي \omterm{def:g10:lewis-and-shape:multiple-bond}{رابطة ثلاثية}
وزوجًا حرًا واحدًا على كل \omterm{def:g7:atoms-and-molecules:atom}{ذرة}:
\chemfig{\charge{180=\:}{N}~\charge{0=\:}{N}}.
\end{solution}

\begin{solution}{exo:g10:lewis-and-shape:8}
\chemfig{H-C(-[2]H)(-[6]H)-\charge{90=\:,270=\:}{O}-H}: أربع روابط حول الكربون،
رباعي الأوجه؛ وحول الأكسجين رابطتان وزوجان حران، منحنٍ.
\end{solution}

\begin{solution}{exo:g10:lewis-and-shape:9}
في \ce{CO2} للكربون مجموعتان (رابطتان ثنائيتان) ولا \omterm{def:g10:lewis-and-shape:lone-pair}{زوج حر}: فتتجهان في
اتجاهين متعاكسين. وفي \ce{H2O} للأكسجين أربع مجموعات (رابطتان، وزوجان حران)
في ترتيب رباعي الأوجه: فتصنع الرابطتان
زاوية.
\end{solution}

\begin{solution}{exo:g10:lewis-and-shape:10}
\chemfig{C(-[3]H)(-[5]H)=C(-[1]H)-[7]H}: حول كل كربون ثلاث مجموعات
(\omterm{def:g10:lewis-and-shape:multiple-bond}{رابطة ثنائية}، ورابطتان أحاديتان): مثلث مستوٍ، وزواياه قريبة من
\ang{120}.
\end{solution}

\begin{solution}{exo:g10:lewis-and-shape:11}
الكربون في المركز، بأربع روابط أحادية إلى \omterm{def:g7:atoms-and-molecules:atom}{ذرات} كلور لكل منها ثلاثة أزواج حرة:
أربع مجموعات حول الكربون، أي رباعي أوجه.
\end{solution}

\begin{solution}{exo:g10:lewis-and-shape:12}
الإيثانول، \chemfig{H-C(-[2]H)(-[6]H)-C(-[2]H)(-[6]H)-\charge{90=\:,270=\:}{O}-H}،
والميثوكسي ميثان،
\chemfig{H-C(-[2]H)(-[6]H)-\charge{90=\:,270=\:}{O}-C(-[2]H)(-[6]H)-H}.
\end{solution}

\begin{solution}{exo:g10:lewis-and-shape:13}
\omterm{def:g10:lewis-and-shape:lone-pair}{الزوج الحر} يشغل حيزًا أكبر قليلًا من \omterm{def:g10:lewis-and-shape:lone-pair}{الزوج الرابط} ويقرّب الروابط بعضها من بعض.
وللأمونياك \omterm{def:g10:lewis-and-shape:lone-pair}{زوج حر} واحد، وللماء زوجان: فتصغر الزاوية من الميثان إلى الأمونياك
إلى الماء.
\end{solution}

\begin{solution}{exo:g10:lewis-and-shape:14}
\chemfig{H-C~\charge{0=\:}{N}}: مجموعتان حول الكربون، خطي.
\chemfig{H-C(-[6]H)=\charge{45=\:,315=\:}{O}}: ثلاث مجموعات حول الكربون، مثلث
مستوٍ.
\end{solution}

\begin{solution}{exo:g10:lewis-and-shape:15}
$3 \times 6 = 18$ إلكترونًا، أي 9 أزواج:
\chemfig{\charge{180=\:,270=\:}{O}=\charge{90=\:}{O}-\charge{0=\:,90=\:,270=\:}{O}}. وللأكسجين المركزي
ثلاث مجموعات (\omterm{def:g10:lewis-and-shape:multiple-bond}{رابطة ثنائية}، ورابطة أحادية، \omterm{def:g10:lewis-and-shape:lone-pair}{وزوج حر}): \omterm{def:g7:atoms-and-molecules:molecule}{فالجزيء}
منحنٍ.
\end{solution}

\begin{solution}{pb:g10:lewis-and-shape:1}
\textbf{1.} \chemfig{H-C(-[2]H)(-[6]H)-H}, \chemfig{H-\charge{90=\:}{N}(-[6]H)-H},
\chemfig{H-\charge{90=\:,270=\:}{O}-H}.

\textbf{2.} أربع مجموعات في كل منها.

\textbf{3.} لا شيء في الميثان، وواحد في الأمونياك، واثنان في الماء.

\textbf{4.} الكربون يكوّن أربع روابط؛ والأكسجين يكوّن رابطتين (فمجموعتاه
الأخريان زوجان حران، لا يمثلهما الطقم بعصيّ).

\textbf{5.} \chemfig{H-C(-[2]H)(-[6]H)-C(-[2]H)(-[6]H)-H}: عصا واحدة تربط
كرتي الكربون، وتلزم ست كرات هيدروجين.

\textbf{6.} رباعي الأوجه، وهرمي، ومنحنٍ.

\textbf{7.} الأزواج الحرة تشغل حيزًا أكبر من الأزواج الرابطة وتضغط الروابط
بعضها على بعض؛ والماء، بزوجيه الحرين، له أصغر زاوية.

\textbf{8.} ثقوب كرة الكربون تصنع \ang{109.5}؛ والزاوية الحقيقية للماء
\ang{104.5}: فكرة الأكسجين تحتاج إلى ثقبين خاصين بها على \ang{104.5}.

\textbf{9.} \ang{180}: فالمجموعتان حول الكربون تتجهان في اتجاهين متعاكسين،
\omterm{def:g7:atoms-and-molecules:molecule}{والجزيء} خطي.

\textbf{10.} الرأسان اللذان ليسا على الضلع نفسه رأسان متقابلان في وجه واحد:
فالمسافة بينهما قطر الوجه، $\sqrt{a^2 + a^2} = a\sqrt{2}$.

\textbf{11.} المركز في منتصف قطر المكعب:
$a\sqrt{3}/2$.

\textbf{12.} H--H: $\sqrt{2} \approx \qty{1.41}{cm}$؛ C--H:
$\sqrt{3}/2 \approx \qty{0.87}{cm}$.

\textbf{13.} كل هيدروجين في رأس من رؤوس المكعب والكربون في مركزه: وكل رأس
على المسافة نفسها من المركز.

\textbf{14.} المثلث C\,H$_1$\,H$_2$ متساوي الساقين (C\,H$_1$ = C\,H$_2$)؛
والمستقيم الواصل من رأسه إلى منتصف القاعدة عمودي على
القاعدة.

\textbf{15.} $\sin \widehat{\text{M\,C\,H}_1} = \dfrac{\text{M\,H}_1}{\text{C\,H}_1}
= \dfrac{a\sqrt{2}/2}{a\sqrt{3}/2} = \sqrt{2/3} \approx 0.816$.

\textbf{16.} نصف الزاوية $\approx \ang{54.7}$؛ والزاوية H--C--H $\approx
2 \times 54.7 = \ang{109.5}$.

\textbf{17.} إنها زاوية رباعي الأوجه \ang{109.5} الواردة في الجدول.

\textbf{18.} يجب أن تُثقب الثقوب على زاوية \ang{109.5} بعضها من بعض.
\end{solution}
